% TEMPLATE-MILC-v3.tex - plantilla maestra UNICA de las guias MILC v3.
%
% La usan las cuatro familias de guias, cada una con su driver:
%   · Grados 6-11          -> scripts/build-guias-g11.py
%   · Bebras               -> scripts/build-guias-bebras.py
%   · Semillero            -> scripts/build-guias-semillero.py
%   · Territorio interior  -> scripts/build-guias-territorio-interior.py
%
% Antes habia tres copias casi identicas (~400 lineas iguales, 6 distintas):
% cada mejora habia que aplicarla tres veces, y en la practica no se hacia
% (el motor de recursos vivia solo en dos de ellas). Lo que varia entre
% familias esta parametrizado en 6 placeholders que cada driver llena
% (se nombran aqui SIN su sintaxis real: el builder reemplaza por texto plano,
% tambien dentro de los comentarios, y un valor multilinea rompe el preambulo):
%
%   HEADER_IZQ        encabezado izquierdo de pagina
%   HEADER_DER        encabezado derecho de pagina
%   PORTADA_ETIQUETA  rotulo sobre el numero grande (GRADO/MOMENTO/MODULO)
%   PORTADA_SERIE     linea de serie de la portada
%   PORTADA_ID        identificador de la guia en la portada
%   PORTADA_CIRCULO   el circulito de grado (°), o vacio si no aplica
%   PDF_TITULO        titulo en texto plano (sin \textbf ni comillas TeX) para
%                     los metadatos del PDF (/Title); lo produce lib_xelatex.texto_pdf
%
% Composicion (auditoria 2026-09-03): las bandas de estacion piden espacio
% (\Needspace*) y prohiben el salto justo despues (\nopagebreak), asi no quedan
% huerfanas al pie; las cajas largas (softbox, drawbox, infoband, puentebox) son
% `breakable`, asi una caja de media pagina ya no arrastra todo a la siguiente
% dejando la anterior al 40 %. Todo texto sobre mostaza o verde va en milcNegro
% (blanco sobre #E5B400 da 1,93:1 y sobre #58A923 2,95:1; ninguno pasa WCAG AA).
% El resto de claves las documenta content/guias/_SCHEMA.md. El script valida
% que no queden placeholders sin reemplazar antes de escribir el archivo final.

% Etiquetado PDF (latex-lab): /Lang, estructura y texto alternativo de las
% figuras, para lectores de pantalla. Debe ir ANTES de \documentclass.
\DocumentMetadata{lang=es-CO,tagging=on}
\documentclass[11pt,letterpaper]{article}
% headheight=24pt: la cabecera derecha lleva dos lineas (identidad de la guia
% y estacion actual). headsep se acorta en la misma medida para que el bloque
% de texto quede exactamente donde estaba con headheight=12pt.
\usepackage[margin=2.0cm,headheight=24pt,headsep=13pt]{geometry}
\usepackage[table]{xcolor}
\usepackage{fontspec}
\usepackage[spanish]{babel}
\usepackage{tikz}
% Librerías TikZ para los diagramas de las guías (bloque `recursos:`):
%  · babel  → babel en español vuelve activos caracteres como «>», y TikZ los usa
%             en sus opciones (>=stealth, ->). Sin esta librería, cualquier
%             diagrama con flechas revienta con «\language@active@arg> has an
%             extra }». Debe ir ANTES de que se dibuje nada.
%  · positioning        → sintaxis «below left=2pt and 3pt».
%  · decorations.pathreplacing → llaves ({brace}) para señalar rangos.
\usetikzlibrary{babel,positioning,decorations.pathreplacing}
\usepackage{graphicx}
% Circuitos: el trabajo de electrónica se hace hoy con micro:bit y sus sensores
% integrados (MakeCode, sin cableado), así que no se necesitan esquemáticos.
% Si algún día entran montajes con protoboard, basta descomentar esta línea:
% los diagramas .tex/.tikz declarados en `recursos:` ya se incrustan solos.
% \usepackage{circuitikz}
\usepackage[most]{tcolorbox}
\usepackage{tabularx}
\usepackage{array}
\usepackage{fancyhdr}
\usepackage{titlesec}
\usepackage[document]{ragged2e}  % bandera a la derecha en todo el cuerpo
\usepackage{enumitem}
\usepackage{microtype}
\usepackage{etoolbox}
\usepackage{needspace}
% hyperref al final del preambulo: metadatos del PDF (titulo, autor, idioma,
% familia), marcadores por estacion (\pdfbookmark) y enlaces sin recuadro.
\usepackage[hidelinks,
  pdftitle={Imagen y texto — el ritmo de la lectura},
  pdfauthor={PhD. Álvaro Cárdenas Orozco},
  pdfsubject={Tecnología e Informática · Grado 9 · Guía 15 · MILC v3},
  pdflang={es-CO},
  bookmarksopen=true,bookmarksnumbered=false]{hyperref}

% Helvetica Neue no trae la flecha U+2192, asi que cada «→» escrito literalmente
% en el YAML se compilaba a NADA, en silencio (462 flechas en 61 guias: rutas del
% tipo «paleta → tipografia → lockup» salian rotas). Mapeamos el caracter a su
% version matematica, que si tiene glifo. Asi el autor puede seguir escribiendo
% «→» comodamente en el YAML.
\usepackage{newunicodechar}
\newunicodechar{→}{\ensuremath{\rightarrow}}
% Mismo problema con los simbolos de marca y vinetas: el «✓» de «una palomita ✓»
% salia como cuadrito vacio en el PDF de todas las guias que lo usaban.
\usepackage{pifont}
\newunicodechar{✓}{\ding{51}}
\newunicodechar{✗}{\ding{55}}
\newunicodechar{✔}{\ding{52}}
\newunicodechar{•}{\textbullet}
\newunicodechar{≥}{\ensuremath{\geq}}
\newunicodechar{≤}{\ensuremath{\leq}}

\setmainfont{Helvetica Neue}
\setsansfont{Helvetica Neue}
\setlength{\parindent}{0pt}
\setlength{\parskip}{6pt}
% Sin particion de palabras: en 6.º-7.º una palabra cortada por guion al final
% de linea («Comuni-cacion») frena la lectura mucho mas que una linea corta.
\hyphenpenalty=10000
\exhyphenpenalty=10000
\pretolerance=2000
\tolerance=3000
\setlength{\emergencystretch}{3em}
\linespread{1.10}
\renewcommand{\arraystretch}{1.25}
\setlist{leftmargin=5mm,itemsep=1mm,topsep=1mm}
\newcolumntype{Y}{>{\RaggedRight\arraybackslash}X}

\definecolor{milcMagenta}{HTML}{E6007E}
\definecolor{milcVino}{HTML}{5A0038}
\definecolor{milcTurquesa}{HTML}{0093A5}
\definecolor{milcVerde}{HTML}{58A923}
\definecolor{milcMostaza}{HTML}{E5B400}
\definecolor{milcOcre}{HTML}{B86600}
\definecolor{milcNegro}{HTML}{241816}
\definecolor{milcGris}{HTML}{F7F7F4}
\definecolor{milcLinea}{HTML}{E8E2DD}
\definecolor{milcGradePrimary}{HTML}{7C3AED}
\definecolor{milcGradeDark}{HTML}{4C1D95}
\definecolor{milcGradeSoft}{HTML}{EDE9FE}
\definecolor{milcGradeAccent}{HTML}{CFE7FF}
\definecolor{milcGradeText}{HTML}{FFFFFF}
\color{milcNegro}

\pagestyle{fancy}
\fancyhf{}
% Cabecera de dos lineas: identidad de la guia arriba y, a la derecha y en
% gris, la estacion en curso (\markright desde \milcestacion). Antes de la
% primera estacion la segunda linea va vacia; el \strut mantiene la altura
% para que la linea de arriba no baile entre paginas.
\lhead{\small Tecnología e Informática · Grado 9\\[-1pt]{\footnotesize\strut}}
\rhead{\small Guía 15 · MILC v3\\[-1pt]{\footnotesize\strut\color{milcNegro!65}\rightmark}}
\cfoot{\small\thepage}
\renewcommand{\headrulewidth}{0pt}

% Sobre mostaza (#E5B400) y verde (#58A923) el texto va en milcNegro; sobre el
% resto de la paleta, en blanco. Se usa dentro de fonttitle/fontupper, donde un
% \color posterior manda sobre coltitle.
\newcommand{\milctintasobre}[1]{%
  \ifboolexpr{test{\ifstrequal{#1}{milcMostaza}} or test{\ifstrequal{#1}{milcVerde}}}%
    {\color{milcNegro}}{\color{white}}}

\newtcolorbox{titlebox}[1]{
  enhanced,colback=#1,colframe=#1,arc=7mm,boxrule=0pt,
  left=7mm,right=7mm,top=3mm,bottom=3mm,
  before skip=8pt,after skip=8pt,
  fontupper=\sffamily\bfseries\Large\color{white}
}

\newtcolorbox{softbox}[3]{
  enhanced,breakable,pad at break=3mm,
  colback=#2,colframe=#1,borderline west={4pt}{0pt}{#1},
  arc=2mm,boxrule=.4pt,left=4mm,right=4mm,top=3mm,bottom=3mm,
  before skip=6pt,after skip=8pt,title={#3},coltitle=white,
  colbacktitle=#1,fonttitle=\sffamily\bfseries\milctintasobre{#1},fontupper=\RaggedRight,
  attach boxed title to top left={xshift=3mm,yshift=-2mm},
  boxed title style={arc=2mm, boxrule=0pt}
}

\newtcolorbox{drawbox}[2]{
  enhanced,breakable,pad at break=4mm,
  colback=white,colframe=#1,arc=4mm,boxrule=1pt,
  left=5mm,right=5mm,top=4mm,bottom=4mm,before skip=8pt,after skip=8pt,
  title={#2},coltitle=white,colbacktitle=#1,fonttitle=\sffamily\bfseries\milctintasobre{#1},
  fontupper=\RaggedRight,
  attach boxed title to top left={xshift=4mm,yshift=-2mm},
  boxed title style={arc=3mm, boxrule=0pt}
}

\newtcolorbox{infoband}[1]{
  enhanced,breakable,pad at break=2mm,colback=#1!8,colframe=#1!50,arc=2mm,boxrule=.5pt,
  left=4mm,right=4mm,top=2mm,bottom=2mm,before skip=4pt,after skip=4pt,
  fontupper=\small\RaggedRight
}

% Puente narrativo entre secciones (contrato editorial Conéctate).
% Da continuidad: "Acabas de X. Ahora vas a Y porque Z."
% Visualmente: banda izquierda en color del grado, texto en itálica pequeña.
\newtcolorbox{puentebox}{
  enhanced,breakable,pad at break=2mm,colback=milcGradeSoft,colframe=milcGradeDark,
  borderline west={3pt}{0pt}{milcGradeDark},
  arc=0mm,boxrule=0pt,left=5mm,right=4mm,top=2.5mm,bottom=2.5mm,
  before skip=4pt,after skip=8pt,
  fontupper=\itshape\small\color{milcNegro}\RaggedRight
}
% La flecha va en modo matematico a proposito: Helvetica Neue NO tiene el glifo
% U+2192, asi que \textrightarrow se compilaba a NADA («Missing character: There
% is no → in font Helvetica Neue») y el puente salia sin flecha. En modo
% matematico la flecha viene de la fuente matematica y si se dibuja.
\newcommand{\puente}[1]{%
  \begin{puentebox}%
    \textcolor{milcGradeDark}{$\rightarrow$}\hspace{2mm}#1%
  \end{puentebox}%
}

\newcommand{\linead}[1]{\noindent\color{milcLinea}\rule{#1}{0.5pt}\color{milcNegro}}
\newcommand{\respuesta}[1]{\par\vspace{2mm}\foreach \n in {1,...,#1}{\noindent\color{milcLinea}\rule{\linewidth}{0.5pt}\par\vspace{5mm}}\color{milcNegro}}
\newcommand{\checkbox}{\textcolor{milcGradeDark}{\fbox{\phantom{x}}}\hspace{5pt}}

% ─── Listas aireadas para las guías socioemocionales ────────────────────
% Componer las verificaciones como «\checkbox texto\\» deja el texto que salta
% de línea debajo del cuadrito y apelmaza el bloque. Estos entornos alinean la
% continuación y airean los ítems. Los usa Territorio Interior; cualquier
% familia puede hacerlo.
\newlist{checklista}{itemize}{1}
\setlist[checklista]{
  label=\textcolor{milcGradeDark}{\fbox{\phantom{x}}},
  leftmargin=9mm, labelsep=3.2mm, itemsep=2.4mm,
  topsep=2.5mm, parsep=0pt, partopsep=0pt, before=\RaggedRight
}
\newlist{pasoslista}{enumerate}{1}
\setlist[pasoslista]{
  label=\textcolor{milcGradeDark}{\sffamily\bfseries\arabic*.},
  leftmargin=8mm, labelsep=3mm, itemsep=2.8mm,
  topsep=1mm, parsep=0pt, partopsep=0pt, before=\RaggedRight
}

\newcommand{\phasecard}[4]{
\begin{tcolorbox}[enhanced,colback=#3,colframe=#2,arc=3mm,boxrule=.5pt,height=3.05cm,valign=center,halign=center,left=1.5mm,right=1.5mm,top=2mm,bottom=2mm]
{\sffamily\bfseries\fontsize{22}{24}\selectfont\textcolor{#2}{#1}}\par
{\sffamily\bfseries\small #4}
\end{tcolorbox}
}

% ── Piezas del rediseno 2026 (legibilidad 6.º-11.º) ───────────────────
% Banda de estacion: reemplaza el titlebox macizo. Titulo a la izquierda,
% dato practico (tiempo/modalidad) a la derecha, en una sola linea.
% Nunca huerfana: pide 9 lineas libres (o salta de pagina rellenando con
% \vfill) y prohibe el salto justo despues de la banda. Nueve y no seis:
% la caja partible que sigue necesita ~80pt (titulo adosado, relleno y dos
% lineas) para arrancar en la misma pagina; con menos, tcolorbox la manda
% entera a la siguiente y la banda se queda sola. El penalty 10000 va
% en `after`, ANTES del espacio posterior: un `after skip` (aunque sea 0pt)
% es un \vskip tras la caja, y ese glue es un punto de corte legal que un
% \nopagebreak escrito despues de \end{tcolorbox} ya no cierra. Cada banda
% deja marcador PDF y actualiza la cabecera derecha.
\newcounter{milcestacion}
\newcommand{\milcestacion}[2]{%
\Needspace*{9\baselineskip}%
\stepcounter{milcestacion}%
\pdfbookmark[1]{#2}{milcestacion\themilcestacion}%
\markright{#2}%
\begin{tcolorbox}[enhanced,colback=#1,colframe=#1,arc=2mm,boxrule=0pt,
  left=5mm,right=5mm,top=2.6mm,bottom=2.6mm,before skip=11pt,
  after={\par\nopagebreak\vskip7pt}]
{\sffamily\bfseries\fontsize{15}{18}\selectfont\milctintasobre{#1}#2}
\end{tcolorbox}}

% Tarjeta de paso numerada: sustituye las tablas «Pilar / Aplicacion», que
% eran formato de consulta docente, no de estudiante. El numero va en un
% circulo sobre el borde para que la secuencia se lea de un vistazo.
\newcommand{\milcpaso}[3]{%
\Needspace{4\baselineskip}%
\begin{tcolorbox}[enhanced,colback=white,colframe=milcLinea,arc=1.5mm,boxrule=.9pt,
  left=14mm,right=4mm,top=3mm,bottom=3mm,before skip=4pt,after skip=4pt,
  fontupper=\RaggedRight,
  overlay={\node[anchor=north west,circle,fill=#2,text=white,inner sep=0pt,
    minimum size=7.4mm,font=\sffamily\bfseries\small]
    at ([xshift=3.6mm,yshift=-3.2mm]frame.north west) {#1};}]
#3
\end{tcolorbox}}

% Casilla marcable de tamano util con lapiz (la anterior era \fbox{\phantom{x}},
% de alto variable segun el texto que la seguia).
\newcommand{\milcmarca}{\framebox[3.8mm]{\rule{0pt}{3.4mm}}}

% Fila marcable de lista suelta (checks de la estacion 1).
\newcommand{\milccheckrow}[1]{%
\par\vspace{1.8mm}\noindent
\makebox[6.4mm][l]{\raisebox{-.55mm}{\textcolor{milcNegro}{\milcmarca}}}%
\parbox[t]{\dimexpr\linewidth-6.4mm\relax}{\RaggedRight #1}\par}

% Celda marcable dentro de la rubrica (tabla de autoevaluacion). Misma
% sangria francesa, pero pensada para vivir dentro de una columna X.
\newcommand{\milcopcion}[1]{%
\makebox[5.2mm][l]{\raisebox{-.55mm}{\textcolor{milcNegro}{\milcmarca}}}%
\parbox[t]{\dimexpr\linewidth-5.2mm\relax}{\RaggedRight #1}}

% ── Recursos visuales de la guía ──────────────────────────────────────
% Declarados en el bloque `recursos:` del YAML y emitidos por el builder.
% \guiaFigura   → imagen rasterizada o PDF (foto, carta celeste, montaje).
% \guiaDiagrama → fuente TikZ/circuitikz (.tex) incrustada como vector:
%                 es la vía para circuitos y esquemáticos editables.
% [#1] = texto alternativo (alt) · #2 = ruta relativa al .tex · #3 = pie
% El alt llega al PDF etiquetado (/Alt) y lo lee el lector de pantalla.
\newcommand{\guiaFigura}[3][]{%
\begin{center}
\includegraphics[width=0.88\linewidth,height=0.38\textheight,keepaspectratio,alt={#1}]{#2}\\[1.5mm]
{\footnotesize\itshape\textcolor{milcNegro!75}{#3}}
\end{center}
}

\newcommand{\guiaDiagrama}[2]{%
\begin{center}
\input{#1}\\[1.5mm]
{\footnotesize\itshape\textcolor{milcNegro!75}{#2}}
\end{center}
}

\begin{document}
\thispagestyle{empty}

% ── Portada ───────────────────────────────────────────────────────────
% Antes: un rectangulo de color a pagina completa. Se veia bien en pantalla,
% pero en la fotocopiadora del colegio se comia el toner y salia gris sucio.
% Ahora la banda ocupa el tercio superior y el resto va en blanco.
\begin{tikzpicture}[remember picture,overlay]
  \path[fill=milcGradePrimary]
    (current page.north west) rectangle ([yshift=-10.6cm]current page.north east);

  \node[anchor=north west,text=milcGradeText] at ([xshift=2.0cm,yshift=-1.55cm]current page.north west)
    {\sffamily\bfseries\fontsize{12}{14}\selectfont GRADO};
  \node[anchor=north west,text=milcGradeText,opacity=.85] at ([xshift=2.0cm,yshift=-2.35cm]current page.north west)
    {\sffamily\bfseries\fontsize{8.5}{10}\selectfont SERIE GUÍAS · TECNOLOGÍA E INFORMÁTICA};
  \node[anchor=north east,text=milcGradeText,opacity=.85,align=right] at ([xshift=-2.0cm,yshift=-1.55cm]current page.north east)
    {\sffamily\bfseries\fontsize{9}{12}\selectfont I.E. Sor María Juliana\\Cartago, Valle del Cauca};

  \node[anchor=west,text=milcGradeText] (gradonum) at ([xshift=1.85cm,yshift=-7.1cm]current page.north west)
    {\sffamily\bfseries\fontsize{112}{112}\selectfont 9};
  % El circulito de grado (°) solo aplica a las guias de grado («11°»); en
  % Bebras y Semillero el numero grande es un momento/modulo, no un grado.
  % Se posiciona relativo al nodo `gradonum`, no en coordenadas absolutas.
  \draw[milcGradeText,line width=1.6pt]
    ([xshift=2.0mm,yshift=-4.5mm]gradonum.north east) circle (.15cm);

  \node[anchor=west,text=milcGradeText,text width=11.4cm,align=left] at ([xshift=1.5cm,yshift=0.5cm]gradonum.east)
    {
     {\sffamily\bfseries\fontsize{9}{11}\selectfont\MakeUppercase{Período 2 · Diseño editorial digital}}\\[3mm]
     {\sffamily\bfseries\fontsize{21}{25}\selectfont\setlength{\baselineskip}{27pt}Imagen y texto ---\\[4pt]quién habla\\[4pt]primero}\\[3.5mm]
     {\sffamily\bfseries\fontsize{15}{18}\selectfont Guía 15-9-TIC}
    };
\end{tikzpicture}

\vspace*{9.4cm}
\vfill

{\sffamily\bfseries\fontsize{8}{10}\selectfont\textcolor{milcGradeDark}{LO QUE TE LLEVAS HOY}}

\begin{tcolorbox}[enhanced,colback=white,colframe=milcNegro,arc=2mm,boxrule=1.6pt,
  left=5mm,right=5mm,top=4mm,bottom=4mm,before skip=3pt,after skip=12pt,
  fontupper=\RaggedRight\fontsize{11}{15.5}\selectfont]
Dos páginas dobles de tu revista dibujadas a mano en hojas A4 horizontales, con la cuadrícula trazada a lápiz suave, la imagen ubicada, un titular real y el cuerpo simulado. Cada una aplica el par tipográfico y la paleta que ya definiste, y declara qué relación tienen la imagen y el texto: si se refuerzan, si contrastan o si cuentan una secuencia.
\end{tcolorbox}

\begin{tcolorbox}[enhanced,colback=milcGris,colframe=milcGris,arc=2mm,boxrule=0pt,
  left=5mm,right=5mm,top=3.5mm,bottom=3.5mm,after skip=12pt]
{\sffamily\bfseries\fontsize{8}{10}\selectfont\textcolor{milcNegro!55}{ESTA GUÍA ES DE}}\\[3mm]
\begin{tabularx}{\linewidth}{X p{3.2cm} p{3.2cm}}
\linead{\linewidth} & \linead{3.0cm} & \linead{3.0cm}\\[-1mm]
{\fontsize{8.5}{10}\selectfont\textcolor{milcNegro!55}{Nombre completo}} &
{\fontsize{8.5}{10}\selectfont\textcolor{milcNegro!55}{Grupo}} &
{\fontsize{8.5}{10}\selectfont\textcolor{milcNegro!55}{Fecha}}\\
\end{tabularx}
\end{tcolorbox}

\vfill

\noindent\textcolor{milcLinea}{\rule{\linewidth}{1pt}}\\[2mm]
{\sffamily\bfseries\fontsize{9.5}{12}\selectfont Autor: Dr. Álvaro Cárdenas Orozco}\\[1mm]
{\sffamily\fontsize{8.5}{11}\selectfont\textcolor{milcNegro!55}{Metodología MILC · apertura ancestral · pensamiento computacional · triángulo Dussel–estoicismo–Floridi}}

\newpage

\section*{Apertura: Imagen y texto --- el ritmo de la lectura}

\begin{softbox}{milcGradeDark}{milcGradeSoft}{Saber ancestral}
En los pueblos del Cauca hay un lugar donde se aprende a escuchar: la \textbf{tulpa}, el fogón de tres piedras en el centro de la casa. El Tejido de Comunicación de la ACIN recoge la explicación de los mayores kiwe the del resguardo de Toez, en Caloto: «las tres piedras que forman la tulpa son traídas con todos los análisis espirituales de los mayores… es aquí donde está la comunicación propia» (2018). Cuatro comuneros del pueblo kokonuko lo cuentan igual en la revista del CRIC: alrededor del fogón, de noche y después de comer, los abuelos contaban historias a hijos y nietos. Fíjate en el detalle. Nadie llegaba allí a hablar: se llegaba a estar y a oír. La \textbf{cara de exclusión}: la tulpa es también un espacio de trabajo doméstico repartido de forma desigual, y el CRIC usa el nombre para sus espacios políticos. Es institución de gobierno, no una técnica de aula.\par\smallskip{\footnotesize\textcolor{milcNegro!70}{\textbf{Fuente.} Tejido de Comunicación ACIN. (2018, 17 de agosto). \emph{Alrededor de la tulpa como principio de la comunicación ancestral, inició el diálogo con los pares académicos de la UAIIN-CRIC}. Çxhab Wala Kiwe.}}
\end{softbox}

\begin{softbox}{milcTurquesa}{milcGris}{Saber contemporáneo}
En una página doble, la imagen y el texto tienen que turnarse. Si los dos gritan a la vez, el lector no oye a ninguno. Hay tres relaciones posibles y conviene elegir una a propósito. \textbf{Refuerzo}: los dos dicen lo mismo por canales distintos, y sirve cuando el mensaje debe quedar claro sin margen de duda. \textbf{Contraste}: la imagen dice algo que el texto matiza o contradice, y sirve para hacer pensar. \textbf{Secuencia}: la imagen empieza la historia y el texto la continúa, o al revés. Y hay un patrón que ordena casi cualquier composición: la imagen captura y el texto profundiza. El ojo entra por la imagen porque no necesita descifrarla, y la mente sigue por el texto. Por eso la pregunta que resuelve un spread es siempre la misma: qué quiero que vea primero, qué después y qué no necesita mirar.
\end{softbox}

\begin{softbox}{milcMostaza}{milcGris}{Pregunta-puente}
\textit{A la tulpa se llega a estar y a oír, no a hablar encima del otro. En tu página doble, ¿quién habla primero y quién espera su turno?}
\end{softbox}

\begin{softbox}{milcVerde}{milcGris}{Saber y saber hacer}
Al terminar podrás: (1) \textbf{identificar} qué relación tienen imagen y texto en páginas dobles reales; (2) \textbf{explicar} las tres relaciones y el patrón de la imagen que captura y el texto que profundiza; (3) \textbf{crear} dos páginas dobles a mano con tu sistema visual y la relación declarada en cada una.
\end{softbox}



\small
\begin{tabularx}{\textwidth}{>{\bfseries}p{3.0cm}Y}
\rowcolor{milcGradePrimary!12}Autor & Dr. Álvaro Cárdenas Orozco\\
DBA & Producción de piezas editoriales digitales con criterio estético y ético (MEN, T\&I)\\
Referentes & Tipografía colombiana · Diseño centrado en accesibilidad · Floridi\\
Producto final & Dos páginas dobles de tu revista dibujadas a mano en hojas A4 horizontales, con la cuadrícula trazada a lápiz suave, la imagen ubicada, un titular real y el cuerpo simulado. Cada una aplica el par tipográfico y la paleta que ya definiste, y declara qué relación tienen la imagen y el texto: si se refuerzan, si contrastan o si cuentan una secuencia.\\
\end{tabularx}
\normalsize

\puente{Ya tienes cuadrícula, tipografía y paleta. Hoy los pones a trabajar juntos en una página doble, que es la unidad real de una revista. En la sesión 6 llevas todo esto a una herramienta digital.}

\milcestacion{milcGradeDark}{Ruta MILC: así vamos a aprender}
\begin{tabularx}{\textwidth}{>{\centering\arraybackslash}X>{\centering\arraybackslash}X>{\centering\arraybackslash}X>{\centering\arraybackslash}X}
\phasecard{1}{milcVerde}{milcGris}{Escucha\\[-1mm]\small Observo, escucho y pregunto.} &
\phasecard{2}{milcTurquesa}{milcGris}{Entender\\[-1mm]\small Sistematizo y ordeno.} &
\phasecard{3}{milcMagenta}{milcGris}{Hacer\\[-1mm]\small Construyo y aplico.} &
\phasecard{4}{milcVino}{milcGris}{Cierre\\[-1mm]\small Evalúo y reflexiono.}
\end{tabularx}


\puente{Antes de la teoría vas a mirar tres páginas dobles reales y a fijarte en qué miraste primero. Esa primera mirada es la que el diseñador quiso o no consiguió.}

\milcestacion{milcVerde}{Estación 1 de 3 · Escucha}

\begin{softbox}{milcVerde}{milcGris}{Escena para escuchar}
\textbf{Actividad 1 · IDENTIFICA --- Tres páginas dobles reales} (15 min · individual). (1) Busca tres páginas dobles: una de periódico, una de revista y una de un libro ilustrado o un catálogo. (2) Mira cada una durante cinco segundos y anota qué miraste primero. (3) Anota qué miraste después. (4) Escribe si la imagen y el texto dicen lo mismo, se contradicen o se continúan. (5) Marca en cuál compiten entre sí y explica en una línea cómo lo notaste.
\end{softbox}

{\sffamily\bfseries\fontsize{8}{10}\selectfont\textcolor{milcNegro!55}{MARCA CUANDO LO HAYAS HECHO}}
\milccheckrow{Tengo las tres con qué miré primero y qué después.}
\milccheckrow{Escribí qué relación tienen imagen y texto en cada una.}
\milccheckrow{Marqué en cuál compiten y expliqué cómo lo noté.}

\begin{infoband}{milcVerde}


\textbf{Lo que va al cuaderno (Actividad 1 · IDENTIFICA).} \textbf{Título:} «Actividad 1 --- Tres páginas dobles reales». \textbf{Formato:} tabla de 3 filas y 4 columnas (pieza / qué miré primero / qué después / relación entre imagen y texto), con marca en la que compiten. \textbf{Extensión:} un tercio de página. \textbf{Sabes que terminaste cuando} las tres filas tienen nombrada su relación.
\end{infoband}

\puente{Ya lo viste. Ahora vas a nombrar las tres relaciones posibles y a decidir cuál usa cada spread tuyo.}

\milcestacion{milcTurquesa}{Estación 2 de 3 · Entender}

\begin{softbox}{milcTurquesa}{milcGris}{Lo que vas a entender}
\textbf{Actividad 2 · EXPLICA --- Las tres relaciones} (20 min · parejas). (1) Con tu pareja, escriban las tres relaciones con una frase propia cada una. (2) Clasifiquen sus seis piezas. (3) Tomen la que compite y escriban qué cambiarían para que uno de los dos ceda el turno. (4) Escriban en dos líneas para qué sirve cada relación, es decir, cuándo conviene usarla.
\end{softbox}

\milcpaso{1}{milcTurquesa}{\textbf{Las tres relaciones.} Refuerzo: los dos dicen lo mismo por canales distintos. Contraste: la imagen dice algo que el texto matiza. Secuencia: uno empieza la historia y el otro la continúa.}
\milcpaso{2}{milcTurquesa}{\textbf{La imagen captura, el texto profundiza.} El ojo entra por la imagen porque no hay que descifrarla, y la mente sigue por el texto. Cuando el titular es tan grande que compite con la foto, el lector no sabe por dónde empezar.}
\milcpaso{3}{milcTurquesa}{\textbf{Si los dos dicen exactamente lo mismo, uno sobra.} El refuerzo sirve cuando el mensaje no admite dudas, pero repetir por repetir desperdicia media página. La foto debería aportar lo que el texto no puede decir.}
\milcpaso{4}{milcTurquesa}{\textbf{Turnarse no es un adorno.} A la tulpa se llega a estar y a oír. Una página doble donde todo grita a la vez es una conversación donde nadie escucha, y el lector se va sin haber entendido nada.}

\begin{drawbox}{milcTurquesa}{Imagen y texto, en una mirada}
\textbf{Refuerzo:} lo mismo por dos canales. \\[1mm] \textbf{Contraste:} la imagen dice, el texto matiza. \\[1mm] \textbf{Secuencia:} uno empieza, el otro sigue. \\[1mm] \textbf{Patrón:} la imagen captura, el texto profundiza. \\[1mm] \textbf{Pregunta:} ¿qué se mira primero y qué después?
\end{drawbox}

\begin{softbox}{milcMostaza}{milcGris}{Errores comunes y su corrección}
(1) \textbf{Titular e imagen del mismo peso}: compiten y el lector no sabe por dónde entrar. (2) \textbf{Foto que repite el titular}: se pierde media página diciendo dos veces lo mismo. (3) \textbf{Texto pegado al borde de la imagen}: sin aire entre los dos, se leen como una sola mancha. (4) \textbf{No declarar la relación}: si tú no sabes cuál usaste, el lector tampoco. (5) \textbf{Llenar los dos lados por igual}: una página doble no tiene que ser simétrica para estar equilibrada.
\end{softbox}

\begin{infoband}{milcTurquesa}


\textbf{Lo que va al cuaderno (Actividad 2 · EXPLICA).} \textbf{Título:} «Actividad 2 --- Las tres relaciones». \textbf{Formato:} las tres con frase propia, las seis piezas clasificadas, el arreglo de la que compite y las dos líneas sobre cuándo conviene cada relación. \textbf{Extensión:} media página. \textbf{Sabes que terminaste cuando} la pieza que competía tiene escrito qué cede el turno.
\end{infoband}

\puente{Tienes el criterio. Toca dibujar dos páginas dobles a mano, con la cuadrícula visible y la relación declarada.}

\milcestacion{milcMagenta}{Estación 3 de 3 · Hacer}

\begin{softbox}{milcMagenta}{milcGris}{Cómo vas a construir tu producto}
\textbf{Actividad 3 · CREA --- Dos páginas dobles} (30 min · individual). (1) Toma dos hojas A4 horizontales y traza en cada una la cuadrícula a lápiz suave. (2) Decide para cada spread qué historia cuenta y qué relación tendrán imagen y texto. (3) Ubica la imagen y escribe el titular real con tu tipografía display. (4) Simula el cuerpo con líneas, respetando la mancha y el aire. (5) Aplica la paleta señalando dónde va cada color. (6) Enséñaselos a un compañero cinco segundos cada uno y anota qué dice que miró primero.
\end{softbox}

\milcpaso{1}{milcMagenta}{\textbf{Tu entrega son dos spreads y una prueba.} Los dos dibujados con cuadrícula visible, relación declarada y sistema visual aplicado. Y lo que dijo tu compañero después de cinco segundos.}
\milcpaso{2}{milcMagenta}{\textbf{Las dos relaciones tienen que ser distintas.} Dos spreads de refuerzo son el mismo ejercicio dos veces. Prueba uno de refuerzo y otro de contraste o secuencia, que es donde se aprende.}
\milcpaso{3}{milcMagenta}{\textbf{El titular va escrito de verdad.} El cuerpo se puede simular con líneas, pero el titular no: es la pieza que decide el peso visual y hay que verlo con sus palabras reales y su tipografía.}
\milcpaso{4}{milcMagenta}{\textbf{Cinco segundos bastan para la prueba.} Nadie mira una página doble más tiempo antes de decidir si sigue. Si en cinco segundos tu compañero mira lo que no querías, hay que mover algo.}

\begin{drawbox}{milcMagenta}{Antes de entregar}
\checkbox Dos hojas A4 horizontales con la cuadrícula a lápiz suave. \\[1mm] \checkbox Imagen ubicada y titular escrito con la tipografía display. \\[1mm] \checkbox Cuerpo simulado respetando la mancha y el aire. \\[1mm] \checkbox Paleta señalada, con la función de cada color. \\[1mm] \checkbox Relación declarada en cada spread, y distintas entre sí. \\[1mm] \checkbox Un compañero dijo qué miró primero en cada uno.
\end{drawbox}

\begin{softbox}{milcMagenta}{milcGris}{Plantilla de trabajo}
\textbf{Historia.} \textit{Qué cuenta este spread.} \\[1mm] \textbf{Relación.} \textit{Refuerzo, contraste o secuencia.} \\[1mm] \textbf{Entrada.} \textit{Qué quiero que se mire primero.} \\[1mm] \textbf{Sistema.} \textit{Cuadrícula, tipografía y colores aplicados.} \\[1mm] \textbf{Prueba.} \textit{Qué miró primero mi compañero.}
\end{softbox}

\begin{infoband}{milcMagenta}


\textbf{Lo que va al cuaderno (Actividad 3 · CREA).} \textbf{Título:} «Actividad 3 --- Dos páginas dobles». \textbf{Formato:} los dos bocetos con su relación declarada y la nota de qué miró primero tu compañero en cada uno. \textbf{Extensión:} media página. \textbf{Sabes que terminaste cuando} en al menos uno coincidió lo que miró con lo que querías.
\end{infoband}

\puente{Lo que entregas hoy: dos páginas dobles dibujadas. La prueba de calidad: un compañero dice qué miró primero y coincide con lo que querías.}

\milcestacion{milcMostaza}{Producto de la sesión}
\textbf{Dos páginas dobles dibujadas con relación imagen-texto declarada}.
\begin{softbox}{milcMostaza}{milcGris}{Criterios de calidad}
(1) Los dos spreads están en A4 horizontal con la cuadrícula visible a lápiz. (2) La imagen está ubicada y el titular escrito con la tipografía display. (3) Se aplican el par tipográfico y la paleta de las sesiones anteriores. (4) Cada spread declara su relación y las dos son distintas entre sí. (5) Un compañero dijo qué miró primero y coincidió al menos en uno.
\end{softbox}

\puente{Antes de cerrar, mira lo que hiciste desde cinco lados. Quien diseña decide el turno de palabra entre una foto y un párrafo, y casi nunca lo dice en voz alta.}

\milcestacion{milcVino}{Evaluación liberadora · cinco dimensiones}

\begin{softbox}{milcVino}{milcGris}{Sentido de la evaluación}
Evaluar no es solo revisar una nota. Es reconocer qué aprendí, qué cambió en mí, qué emociones manejé, cómo me relaciono con la ciudad y la comunidad, y qué le dejaré a quien viene detrás.
\end{softbox}

\textbf{Mi reflexión liberadora en cinco dimensiones.} Completa cada frase iniciada:

\small
\begin{tabularx}{\textwidth}{>{\bfseries\RaggedRight\arraybackslash}p{3.4cm}Y>{\RaggedRight\arraybackslash}p{4.6cm}}
\rowcolor{milcVino}\color{white}Dimensión & \color{white}\bfseries Mi respuesta (frame starter) & \color{white}\bfseries Algo que descubrí\\
1. Personal & Lo que más cambió en mí fue \linead{6.5cm} & \linead{4.2cm}\\
2. Emocional & La emoción más fuerte fue \linead{2.5cm} y la regulé \linead{3cm} & \linead{4.2cm}\\
3. Ciudadana & En democracia esto importa porque \linead{6cm} & \linead{4.2cm}\\
4. Local (Cartago) & En mi barrio sería útil para \linead{6.5cm} & \linead{4.2cm}\\
5. Intergeneracional & A mi abuela/abuelo le diría \linead{6.5cm} & \linead{4.2cm}\\
\end{tabularx}
\normalsize

\begin{infoband}{milcVino}
\textbf{Para la próxima sustentación.} Vuelve a esta tabla en seis meses. Verás que las respuestas cambian — no porque lo aprendido cambie, sino porque tú cambias. La evaluación liberadora es una conversación contigo a través del tiempo.
\end{infoband}


\puente{Para terminar, tres citas verificadas sobre escuchar y esperar el turno: Dussel sobre el silencio del que todo tiene que escuchar, Marco Aurelio sobre atender a lo que dice el otro, Floridi sobre la atención como bien escaso.}

\milcestacion{milcGradeDark}{Triángulo de pensamiento · cierre filosófico}

\begin{softbox}{milcVino}{milcGris}{Dussel · lente del nosotros}
\textit{«El respeto es silencio, pero no silencio del que nada tiene que decir, sino del que todo tiene que escuchar porque nada sabe del otro como otro.»}

{\footnotesize\itshape\color{milcNegro!70}--- Enrique Dussel · Filosofía de la liberación (1977), §2.6.2.3}

En una página doble, el silencio es el espacio en blanco y el elemento que cede el turno. Que el texto calle un momento para que la imagen hable no es debilidad del texto: es lo que permite que después se le oiga.

\textbf{¿Qué elemento de mi spread no está dispuesto a callarse ni un segundo?}
\end{softbox}
\linead{\linewidth}\\[1mm]
\linead{\linewidth}

\begin{softbox}{milcMostaza}{milcGris}{Estoicismo · Marco Aurelio · Meditaciones VI, 53 (c. 175 d.C.) · cuidado interior}
\textit{«Acostúmbrate a estar con atención a lo que dice el otro, y en cuanto te sea posible intérnate dentro del alma del que hablare contigo.»}

{\footnotesize\itshape\color{milcNegro!70}--- Marco Aurelio · Meditaciones VI, 53 (c. 175 d.C.)}

Diseñar una página doble es imaginar a alguien que todavía no la ha visto. Los cinco segundos de la prueba sirven para eso: para dejar de mirar tu spread como quien lo hizo y mirarlo como quien llega.

\textbf{¿Estoy mirando mi página como el que la hizo o como el que va a llegar a ella?}
\end{softbox}
\linead{\linewidth}\\[1mm]
\linead{\linewidth}

\begin{softbox}{milcTurquesa}{milcGris}{Floridi · lente de la infoesfera}
\textit{«Afirmamos que las capacidades atencionales son un bien finito, precioso y escaso. (trad. propia)»}

{\footnotesize\itshape\color{milcNegro!70}--- The Onlife Initiative (ed. Luciano Floridi) · The Onlife Manifesto (2015), § 4.6}

Por eso el turno importa. Si la imagen y el titular piden atención a la vez, el lector gasta la suya decidiendo dónde mirar en vez de entendiendo. Ordenar la entrada es devolverle ese esfuerzo.

\textbf{¿Cuánta atención le cuesta a alguien decidir por dónde entrar a mi página?}
\end{softbox}
\linead{\linewidth}\\[1mm]
\linead{\linewidth}

\begin{infoband}{milcNegro}
\textbf{Nota para el docente.} Las tres son citas textuales y llevan su referencia debajo. Puedes remitir a la obra en clase.
\end{infoband}

\milcestacion{milcVino}{Mi autoevaluación · marca una casilla por fila}

{\small
\setlength{\extrarowheight}{2.2pt}
\begin{tabularx}{\textwidth}{>{\RaggedRight\arraybackslash\bfseries}p{3.0cm}YYY}
\rowcolor{milcVino}\color{white}\bfseries Criterio & \color{white}\bfseries Lo logré & \color{white}\bfseries En proceso & \color{white}\bfseries Necesito apoyo\\[.6mm]
Comprensión del tema & \milcopcion{Explico las ideas con mis palabras.} & \milcopcion{Reconozco las ideas, falta precisión.} & \milcopcion{Necesito repasar lo básico.}\\[.8mm]
\rowcolor{milcGris}Pensamiento computacional & \milcopcion{Usé los pilares al armar mi producto.} & \milcopcion{Usé algunos pilares.} & \milcopcion{Necesito releer la estación 2.}\\[.8mm]
Aplicación y producto & \milcopcion{Mi producto cumple los criterios.} & \milcopcion{Está armado, con ajustes pendientes.} & \milcopcion{Aún está en construcción.}\\[.8mm]
\rowcolor{milcGris}Reflexión 5 dimensiones & \milcopcion{Respondí cada una con honestidad.} & \milcopcion{Respondí algunas con detalle.} & \milcopcion{Mis respuestas son muy generales.}\\[.8mm]
Triángulo de pensamiento & \milcopcion{Conecté las 3 voces con mi trabajo.} & \milcopcion{Conecté al menos una.} & \milcopcion{No las relaciono con mi proceso.}\\[.8mm]
\end{tabularx}}

\vspace{3mm}
\textbf{Hoy entendí que:}\\[1mm]
\linead{\linewidth}\\[1mm]
\linead{\linewidth}

\textbf{Mi compromiso para la próxima sesión es:}\\[1mm]
\linead{\linewidth}\\[1mm]
\linead{\linewidth}

\begin{softbox}{milcMostaza}{milcGris}{Cita de despedida · Marco Aurelio}
\textit{``El obstáculo en el camino se vuelve el camino. Lo que estorba al camino se vuelve camino.''}\\[1mm]
Cada error de hoy, bien observado, se convierte en pieza del camino que pisarás mañana.
\end{softbox}

\small
\begin{tabularx}{\textwidth}{>{\bfseries}p{3.6cm}Y}
\rowcolor{milcGradeDark!12}Nombre completo: & \linead{10cm}\\
Fecha: & \linead{4cm} \quad Curso/grupo: \linead{4cm}\\
Mi firma: & \linead{9cm}\\
\end{tabularx}
\normalsize

\begin{infoband}{milcGradeDark}
\textbf{Para la próxima sesión.} Llega con los dos spreads dibujados y con lo que dijo tu compañero. Mira una página doble cualquiera durante cinco segundos y anota qué recuerdas después. En la sesión 6 llevas todo esto a una herramienta digital.
\end{infoband}

\end{document}
